\documentclass[exercice]{thedocument}%
\startdatakeys
\source{}
\BO{2026}
\cycle{4}
\competencesDuSocle{CA,RA}
\objectifsApprentissage{A1ac,A2bc,A3bc}
\automatismes{A1ed,A1ci}
\enddatakeys
\begin{document}%
{\itshape Dans cet exercice, toute trace de recherche, même incomplète ou non
fructueuse, sera prise en compte dans l'évaluation.}
\par\vskip15pt\noindent-- AUREL : Belle pêche ! Combien de poissons et de coquillages
vas-tu pouvoir vendre au marché~?
\par\noindent-- MATHIAS : En tout, je vais pouvoir vendre au marché \<n1; poissons
et \<n2; coquillages.\vskip15pt\noindent Mathias est un pêcheur professionnel.
Il veut vendre des paniers contenant des coquillages et des poissons. Il
souhaite concevoir le plus grand nombre possible de paniers identiques. Enfin
il voudrait qu'il ne lui reste aucun coquillage et aucun poisson dans son
congélateur.
\begin{enumerate}
    \item Combien de paniers au maximum Mathias pourra-t-il concevoir ? Justifier.
    \item Quelle sera la composition de chaque panier ? Justifier.
\end{enumerate}\vskip10pt\hfill{\itshape Brevet Métropole La Réunion 2019}
\end{document}
